\documentclass[10pt,a4paper]{amsart}
\usepackage[margin=19mm]{geometry}
\usepackage{amsmath,amssymb,mathtools}
\usepackage[scr=rsfs,cal=euler]{mathalfa}
\usepackage{xfrac}
\usepackage[unicode,hidelinks,bookmarks=false]{hyperref}
\newcommand{\ie}{i.e.\ }
\newcommand{\cf}{cf.\ }
\newcommand{\resp}{respectively\ }
\newcommand{\Subsection}{\subsection}
\providecommand{\nobreakdash}{\nobreak\hbox{-}}
\setlength{\emergencystretch}{2em}
\multlinegap0pt
\usepackage{amssymb}
\usepackage{dsfont}
\newtheorem{defn}{Definition}
\newtheorem{lemma}{Lemma}
\newtheorem{theorem}{Theorem}
\title{Symmetric bilinear forms, superalgebras and integer matrix factorization}
\author{Dan Fretwell, Jenny Roberts}
\date{Source: arXiv:2404.17881v2 (CC BY 4.0)}
\begin{document}
\maketitle

\noindent\emph{Source and licence.} Symmetric bilinear forms, superalgebras and integer matrix factorization. Version arXiv:2404.17881v2 (CC BY 4.0), distributed by arXiv under the Creative Commons Attribution 4.0 International licence, http://creativecommons.org/licenses/by/4.0/, as recorded in the arXiv licence metadata for this exact version and shown on the abstract page https://arxiv.org/abs/2404.17881v2. Full licence notice: \texttt{licenses/arxiv-2404.17881-CC-BY-4.0.txt}.\par\smallskip

\noindent\textit{Editorial scope.} Counted unit: the article's theorem thm:superalg (source0.tex lines 221--224). The article's complete original proof of it is delivered with it (source0.tex lines 226--253, one \texttt{proof} environment of 28 lines). No mathematics is written, repaired or re-derived: the statement and the proof are the article's own lines, in the article's own notation, and no retained line is substituted or re-flowed.  Retained uncounted as the source-local context the proof needs: lines 180--183; lines 203--206.  Nothing was omitted from any retained span, so the package carries no omission record.  No retained line contains a citation, so the wrapper carries no bibliography and no bibliography entry.  No retained line contains a figure environment, an image-inclusion command or drawing code, so no image is delivered and the package declares no asset dependency.  Editorial formatting only, in addition: an AMS \texttt{amsart} preamble that restates the article's own declarations and macros used by the retained lines, namely the environments \texttt{defn}, \texttt{lemma}, \texttt{theorem} on separate counters, together with the standard packages those lines need.  The retained environments print in this document as Definition 1, Lemma 1, Theorem 1 in order of appearance, whereas the article prints the same items with the numbering of its own full document.  The complete native source of the article as distributed in the arXiv e-print for this exact version is bundled unchanged as \texttt{sources/158247/source0.tex}.
\begin{defn}
\label{def:bfend}
For $\phi_1, \phi_2\in\text{End}_K(V)$ we define: \[B(\phi_1, \phi_2) = \text{Tr}(\phi_1^\dagger\phi_2),\] where $\phi_1^{\dagger}$ is the adjoint of $\phi_1$ with respect to $B$.
\end{defn}

\begin{lemma}
\label{lem:decomp}
There exists an orthogonal decomposition \[\text{End}_K(V) = E^{(0)}(B,w)\oplus E^{(1)}(B,w)\] with respect to the symmetric bilinear form in Definition \ref{def:bfend}.
\end{lemma}

\begin{theorem}
\label{thm:superalg}
The decomposition in Lemma \ref{lem:decomp} defines a superalgebra structure on $\text{End}_K(V)$, i.e. $E^{(i)}(B,w)E^{(j)}(B,w) \subseteq E^{(i+j \bmod 2)}(B,w)$ for all $i,j\in\{0,1\}$.
\end{theorem}

\begin{proof}
First let $\phi_1,\phi_2\in E^{(0)}(B,w)$ have weights $\lambda_1, \lambda_2\in K$ respectively. Then $\phi_1\phi_2\in E^{(0)}(B,w)$ has weight $\lambda_1\lambda_2$ since:
\begin{align*}
    B(\cdot. (\phi_1\phi_2)(w)) = B(\phi_1^\dagger(\cdot),\phi_2(w)) = \lambda_2 B(\phi_1^\dagger(\cdot), w) = \lambda_1\lambda_2 B(\cdot, w), \\
    B(w, (\phi_1\phi_2)(\cdot)) = B(\phi_1^\dagger(w),\phi_2(\cdot)) = \lambda_1 B(w,\phi_2(\cdot)) = \lambda_1\lambda_2 B(w,\cdot).
\end{align*}
Now let $\phi_3,\phi_4\in E^{(1)}(B,w)$ and suppose that $a_1, b_1, a_2, b_2\in \{w\}_B^{\perp}$ are such that \begin{align*}\phi_3 &= \phi_{B,w,a_1} + \phi_{B,b_1,w},\\ \phi_4 &= \phi_{B,w,a_2} + \phi_{B,b_2,w}.\end{align*} 

From this one checks that $(\phi_3\phi_4)(w) = (\phi_4^\dagger\phi_3^\dagger)(w) = B(a_1,b_2)B(w,w) w$, and so:
\begin{align*}
    B(\cdot, (\phi_3\phi_4)(w)) &= B(a_1,b_2)B(w,w) B(\cdot,w),\\
    B(w,(\phi_3\phi_4)(\cdot)) &= B((\phi_4^\dagger\phi_3^\dagger)(w),\cdot) = B(a_1,b_2)B(w,w)B(w,\cdot)
\end{align*}

Since $B(a_1,b_2)B(w,w)\neq 0$ we see that $\phi_3\phi_4\in E^{(0)}(B,w)$ has weight $B(a_1,b_2)B(w,w)$.

We now show that $\phi_1\phi_3$ and $\phi_3\phi_1$ lie in $E^{(1)}(B,w)$. First note that:
\begin{align*}
    B(w,(\phi_1\phi_3)(w)) &= \lambda_1 B(w,\phi_3(w)) = 0,\\
    B(w,(\phi_3\phi_1)(w)) &= B(\phi_3^\dagger(w),\phi_1(w)) = \lambda_1 B(\phi_3^\dagger(w),w) = 0.
\end{align*}

Finally, if $u,v\in\{w\}_B^{\perp}$, we have:
\begin{align*}
    B(u,(\phi_1\phi_3)(v)) &= B(a_1,v)B(u,\phi_1(w)) = \lambda_1 B(a_1,v)B(u,w) = 0, \\
    B(u,(\phi_3\phi_1)(v)) &= B((\phi_1^\dagger\phi_3^\dagger)(u),v) = B(b_1,u)B(\phi_1^\dagger(w),v) = \lambda_1 B(b_1,u)B(w,v) = 0.
\end{align*}
\end{proof}

\end{document}
